\documentclass[addpoints]{exam}
\usepackage[utf8x]{inputenc}
\usepackage[ngerman]{babel}
\usepackage{listings}
\usepackage[top=1.5cm,bottom=0.5cm,headsep=0.5cm,headheight=3cm,%
left=1.5cm,right=1.5cm]{geometry}

\usepackage[T1]{fontenc}
\usepackage{booktabs} % schöne Tabellen
\usepackage{graphicx}
\usepackage{amsmath,textcomp} %diverses
\usepackage{enumitem}
\usepackage{csquotes}

\lstset{language=Python, tabsize=4, basicstyle=\footnotesize, showstringspaces=false, mathescape=true}
\lstset{literate=%
  {==}{{=\hspace{0.15em}=}}3
  {Ö}{{\"O}}1
  {Ä}{{\"A}}1
  {Ü}{{\"U}}1
  {ß}{{\ss}}1
  {ü}{{\"u}}1
  {ä}{{\"a}}1
  {ö}{{\"o}}1
}
\begin{document}
\pointpoints{Punkt}{Punkte}
\bonuspointpoints{Bonuspunkt}{Bonuspunkte}                     
\renewcommand{\solutiontitle}{\noindent\textbf{Lösung:}%
\enspace}

\newcommand{\pk}[1]{\underline{#1}}            % Unterstreichung für primary key
\newcommand{\fk}[1]{{#1}$\uparrow$}           % Pfeil für foreign key
\newcommand{\pfk}[1]{\underline{#1}$\uparrow$}           % Unterstreichung und Pfeil

\pagestyle{headandfoot}
\runningheadrule

\newcommand{\klaubez}{Aufgaben zu Datenbanken}
\firstpageheader{Informatik }{\klaubez} {\thepage /\numpages}
\runningheader{Informatik }{\klaubez} {\thepage /\numpages}
\newcommand{\pfad}{../aufgabentex}



% Steuerung ueber \mitloesungen (0 = ohne Loesungen, 1 = mit Loesungen).
\providecommand{\mitloesungen}{0}
\ifnum\mitloesungen=1
  \printanswers
\fi


% \frage{datei}: bindet eine Aufgabe ein. Fragetext und Loesungsbox werden in eine Box
% gesetzt und damit nie ueber einen Seitenumbruch getrennt (kein manuelles \newpage noetig).
\newsavebox{\qbox}
\newcommand{\frage}[1]{\setbox\qbox\vbox{\hsize\linewidth\input{\pfad/#1}}\box\qbox}


\begin{questions}


% Aufgaben zu Datenbanken
\frage{cur01}
\frage{cur02}
\frage{cur02a}
\frage{cur03}
\frage{cur03a}
\frage{cur04}
\frage{cur04c}
\frage{cur50}
\frage{cur50a}
\frage{cur50b}
\frage{cur05}
\frage{cur06}
\frage{cur06b}
\frage{cur07}
\frage{cur08}
\frage{cur08b}
\frage{cur09}
\frage{cur09a}
\frage{cur09b}
\frage{cur10a}
\input{\pfad/cur11a}

% -------------------------------------------------
\end{questions}
\end{document}
